\thispagestyle{empty}
\setlength{\parindent}{0in}
\setlength{\parskip}{0em}
\linespread{1.0}\selectfont

\textbf{POLITECHNIKA RZESZOWSKA} \hfill Rzeszów 2026 \\
\textbf{im. Ignacego Łukasiewicza} \\
\textbf{WYDZIAŁ MATEMATYKI I FIZYKI STOSOWANEJ} \\

\begin{center}
	\textbf{STRESZCZENIE  PRACY  DYPLOMOWEJ }
\end{center}

\textbf{Tytuł:} Modelowanie dużych sieci LoRa Mesh \\
\textbf{Autor:} Mikołaj Data \\
\textbf{Promotor:} dr Krzysztof Pupka \\
\textbf{Słowa kluczowe:} LoRa, Meshtastic, sieci kratowe, symulacja, model Longley-Rice, routing
\\

Celem pracy było opracowanie symulatora sieci LoRa Mesh opartych na projekcie Meshtastic. Warstwa logiczna sieci, w tym routing i retransmisje, realizowana jest przez rzeczywiste oprogramowanie \texttt{meshtasticd} uruchamiane w kontenerze Docker dla każdego węzła, a symulowana jest wyłącznie warstwa radiowa. Straty propagacyjne wyznaczane są modelem Longley-Rice na podstawie danych SRTM, uzupełnionym o stratę od budynków, cieniowanie i fading. Odwzorowano również kolizje, półdupleks i wykrywanie zajętości kanału. Model łączy skalibrowano na danych z polskiej sieci Meshtastic (lorastats.pl) i oceniono na niezależnym zbiorze testowym: AUC wyniosło 0,79, a tryb przekazania pakietów odwzorowano w 75\% par węzłów. Trasy w symulacji są średnio o 1,5 skoku krótsze niż w sieci rzeczywistej. Działanie symulatora sprawdzono dla sieci liczących do 80 węzłów. Nadaje się on przede wszystkim do porównywania wariantów konfiguracji sieci.


% Streszczenia PL i EN nie mieszczą się na jednej stronie, więc EN zaczyna się od nowej
\clearpage
\thispagestyle{empty}



\textbf{RZESZOW UNIVERSITY OF TECHNOLOGY} \hfill Rzeszów 2026 \\
\textbf{FACULTY OF MATHEMATICS AND APPLIED PHYSICS } \\

\begin{center}
	\textbf{DIPLOMA THESIS ABSTRACT}
\end{center}

\textbf{Title:} Modelling large LoRa Mesh networks\\
\textbf{Author:} Mikołaj Data \\
\textbf{Supervisor:} Krzysztof Pupka, PhD \\ %% Angielski zapis stopni/tytułów naukowych promotora 

\textbf{Key words:} LoRa, Meshtastic, mesh networks, simulation, Longley-Rice model, routing
\\

The aim of the thesis was to develop a simulator of LoRa Mesh networks based on the Meshtastic project. The logical layer of the network, including routing and retransmissions, is handled by the actual \texttt{meshtasticd} firmware running in a Docker container for each node, while only the radio layer is simulated. Path loss is computed with the Longley-Rice model from SRTM data, extended with building loss, shadowing and fading. Collisions, half-duplex operation and channel activity detection are also modelled. The link model was calibrated on data from the Polish Meshtastic network (lorastats.pl) and evaluated on an independent test set: it achieved an AUC of 0.79, and the forwarding mode was reproduced for 75\% of node pairs. Routes in the simulation are on average 1.5 hops shorter than in the real network. The simulator was tested on networks of up to 80 nodes. It is best suited to comparing network configurations.

\vfill

% Poniżej przykłady angielskich odpowiedników polskich stopni i tytułów naukowych promotora:
% - doktor | dr → Dr Firstname Lastname lub Firstname Lastname, PhD
% - dr inż. → Dr Eng. Firstname Lastname lub Firstname Lastname, PhD Eng.
% - doktor habilitowany | dr hab. → Firstname Lastname, PhD, DSc
% - dr hab. inż. → Firstname Lastname, PhD, DSc Eng.
% - dr hab. prof. PRz → Firstname Lastname, PhD, DSc, Assoc. Prof.
% - prof. dr hab. → Prof. Firstname Lastname
% - prof. dr hab. inż. → Prof. Eng. Firstname Lastname
% - doktor, profesor PRz → Doctor, Firstname Lastname, Associate Professor lub Firstname Lastname, PhD, Assoc. Prof.
